% Options for packages loaded elsewhere
\PassOptionsToPackage{unicode}{hyperref}
\PassOptionsToPackage{hyphens}{url}
\documentclass[
]{article}
\usepackage{xcolor}
\usepackage{amsmath,amssymb}
\setcounter{secnumdepth}{-\maxdimen} % remove section numbering
\usepackage{iftex}
\ifPDFTeX
  \usepackage[T1]{fontenc}
  \usepackage[utf8]{inputenc}
  \usepackage{textcomp} % provide euro and other symbols
\else % if luatex or xetex
  \usepackage{unicode-math} % this also loads fontspec
  \defaultfontfeatures{Scale=MatchLowercase}
  \defaultfontfeatures[\rmfamily]{Ligatures=TeX,Scale=1}
\fi
\usepackage{lmodern}
\ifPDFTeX\else
  % xetex/luatex font selection
\fi
% Use upquote if available, for straight quotes in verbatim environments
\IfFileExists{upquote.sty}{\usepackage{upquote}}{}
\IfFileExists{microtype.sty}{% use microtype if available
  \usepackage[]{microtype}
  \UseMicrotypeSet[protrusion]{basicmath} % disable protrusion for tt fonts
}{}
\makeatletter
\@ifundefined{KOMAClassName}{% if non-KOMA class
  \IfFileExists{parskip.sty}{%
    \usepackage{parskip}
  }{% else
    \setlength{\parindent}{0pt}
    \setlength{\parskip}{6pt plus 2pt minus 1pt}}
}{% if KOMA class
  \KOMAoptions{parskip=half}}
\makeatother
\usepackage{longtable,booktabs,array}
\usepackage{caption}
\captionsetup[table]{skip=6pt}
\newcounter{none} % for unnumbered tables
\usepackage{calc} % for calculating minipage widths
% Correct order of tables after \paragraph or \subparagraph
\usepackage{etoolbox}
\makeatletter
\patchcmd\longtable{\par}{\if@noskipsec\mbox{}\fi\par}{}{}
\makeatother
% Allow footnotes in longtable head/foot
\IfFileExists{footnotehyper.sty}{\usepackage{footnotehyper}}{\usepackage{footnote}}
\makesavenoteenv{longtable}
\usepackage{graphicx}
\makeatletter
\newsavebox\pandoc@box
\newcommand*\pandocbounded[1]{% scales image to fit in text height/width
  \sbox\pandoc@box{#1}%
  \Gscale@div\@tempa{\textheight}{\dimexpr\ht\pandoc@box+\dp\pandoc@box\relax}%
  \Gscale@div\@tempb{\linewidth}{\wd\pandoc@box}%
  \ifdim\@tempb\p@<\@tempa\p@\let\@tempa\@tempb\fi% select the smaller of both
  \ifdim\@tempa\p@<\p@\scalebox{\@tempa}{\usebox\pandoc@box}%
  \else\usebox{\pandoc@box}%
  \fi%
}
% Set default figure placement to htbp
\def\fps@figure{htbp}
\makeatother
\setlength{\emergencystretch}{3em} % prevent overfull lines
\providecommand{\tightlist}{%
  \setlength{\itemsep}{0pt}\setlength{\parskip}{0pt}}
\usepackage{bookmark}
\IfFileExists{xurl.sty}{\usepackage{xurl}}{} % add URL line breaks if available
\urlstyle{same}
% fallback for those not using the hyperref driver hyperxmp:
\makeatletter
\@ifundefined{xmpquote}{\newcommand{\xmpquote}[1]{#1}}{}
\makeatother
\hypersetup{
  hidelinks,
  pdfcreator={LaTeX via pandoc}}

\author{}
\date{}

\begin{document}

\section{Bab 4: Ruang Hidup yang
Terampas}\label{bab-4-ruang-hidup-yang-terampas}

\textbf{CELIOS --- Center of Economic and Law Studies}

\emph{Membedah eskalasi konflik sosial dan perampasan ruang agraria di
balik klaim keberhasilan pembangunan.}

\begin{center}\rule{0.5\linewidth}{0.5pt}\end{center}

\subsection{Metodologi}\label{metodologi}

\textbf{Alur Kausalitas (Ekonomi Politik Ekologi):}
\texttt{Ekspansi\ Industri\ \&\ Proyek\ Strategis} →
\texttt{Perampasan\ Ruang\ Hidup\ \&\ Lahan} →
\texttt{Eskalasi\ Konflik\ Sosial/Agraria}

Tesis dari analisis ini membantah narasi kesejahteraan dengan
memperlihatkan bahwa agresivitas izin konsesi, proyek strategis
nasional, hingga perluasan taman nasional dan pariwisata berbanding
lurus dengan meningkatnya resistensi dan terdepaknya masyarakat lokal
dari ruang kelolanya.

\textbf{Variabel Dampak (Y):} * \textbf{Jumlah Konflik:} Riwayat insiden
letupan konflik agraria historis berdasarkan database independen
masyarakat sipil. * \textbf{Sektor Pemicu:} Tipologi konflik yang
dipecah berdasarkan klasifikasi sektor penyebab dominan.

\textbf{Metode Pengolahan Data:} Analisis menggunakan pendekatan
\emph{Trend Analysis} dan tabulasi silang (\emph{Crosstabulation}).
Menyandingkan matriks kejadian letupan konflik secara sektoral untuk
mengekstraksi fakta episentrum sengketa berdarah.

\begin{center}\rule{0.5\linewidth}{0.5pt}\end{center}

\subsection{Hilirisasi \& Pembangunan Berlumur
Konflik}\label{hilirisasi-pembangunan-berlumur-konflik}

Ekspansi industri ekstraktif dan proyek strategis tidak hanya
menumbangkan daya dukung ekologis, tetapi secara agresif merobek tatanan
kehidupan sosial masyarakat. Data empiris mencatat sejarah panjang
perlawanan akar rumput dengan total terjadinya \textbf{95 letupan
konflik agraria} yang tercatat. Konflik ini bukanlah residu acak
pembangunan, melainkan ekses langsung dari model ekonomi yang sangat
rakus daratan.

Secara mengejutkan, aktor perampas lahan utama tidak hanya didominasi
oleh pertambangan dan perkebunan monokultur, namun meluas ke sekor
\textbf{Kehutanan} (Hutan Lindung, Produksi, Konservasi),
\textbf{Infrastruktur \& PSN} (Bendungan, Transmigrasi, Kawasan
Industri), hingga proyek \textbf{Pariwisata \& Pesisir}. Tiga sektor
utama (Perkebunan, Kehutanan, dan Pertambangan) menyumbang porsi
\textbf{82.1\%} dari keseluruhan catatan konflik. Alih-alih mendapatkan
kucuran kesejahteraan, warga lokal justru seringkali dikriminalisasi,
direpresi, dan diusir dari atas ruang penghidupan historis mereka.

\begin{center}\rule{0.5\linewidth}{0.5pt}\end{center}

\subsection{Ringkasan Metrik Agregat}\label{ringkasan-metrik-agregat}

{\def\LTcaptype{none} % do not increment counter
\begin{longtable}[]{@{}
  >{\raggedright\arraybackslash}p{(\linewidth - 4\tabcolsep) * \real{0.3333}}
  >{\raggedright\arraybackslash}p{(\linewidth - 4\tabcolsep) * \real{0.3333}}
  >{\raggedright\arraybackslash}p{(\linewidth - 4\tabcolsep) * \real{0.3333}}@{}}
\toprule\noalign{}
\begin{minipage}[b]{\linewidth}\raggedright
Indikator
\end{minipage} & \begin{minipage}[b]{\linewidth}\raggedright
Nilai
\end{minipage} & \begin{minipage}[b]{\linewidth}\raggedright
Keterangan
\end{minipage} \\
\midrule\noalign{}
\endhead
\bottomrule\noalign{}
\endlastfoot
\textbf{Total Letupan Konflik} & \textbf{95 kasus} & Insiden perampasan
lahan dan sengketa agraria yang memicu perlawanan sipil. \\
\textbf{Korban Terdampak (Jiwa)} & \textbf{90,582 jiwa} & Jumlah warga
yang kehilangan ruang hidup, digusur, atau terpinggirkan akibat konflik
lahan (bukan korban meninggal). \\
\textbf{Status: Belum Ditangani} & \textbf{51 kasus} & Kasus yang
dibiarkan terkatung-katung tanpa resolusi berkeadilan bagi warga. \\
\textbf{Masyarakat Melawan} & \textbf{59 komunitas} & Kelompok tani dan
masyarakat adat yang berjuang mempertahankan ruang hidup. \\
\textbf{Sektor Perkebunan} & \textbf{25 kasus} & Tumpang tindih Hak Guna
Usaha (HGU) sawit skala masif dengan lahan rakyat. \\
\textbf{Sektor Kehutanan} & \textbf{30 kasus} & Klaim sepihak hutan
produksi dan konservasi yang menggusur masyarakat lokal. \\
\textbf{Sektor Pertambangan} & \textbf{23 kasus} & Operasi pengerukan
lahan dan hilirisasi untuk industri mineral serta nikel. \\
\textbf{Infrastruktur \& PSN} & \textbf{13 kasus} & Penggusuran proyek
strategis nasional seperti bendungan dan jalan. \\
\textbf{Pariwisata \& Pesisir} & \textbf{3 kasus} & Privatisasi pesisir
dan pariwisata super-premium (KEK). \\
\textbf{Keterlibatan Pemerintah} & \textbf{69 kasus} & Andil institusi
negara dan pemerintah daerah dalam sengketa warga. \\
\textbf{Keterlibatan Korporasi} & \textbf{52 kasus} & Perusahaan swasta
asing maupun BUMN yang memonopoli ruang hidup. \\
\end{longtable}
}

\emph{Sumber Analisis Data: Konsorsium Pembaruan Agraria (KPA) / Tanah
Kita}

\begin{center}\rule{0.5\linewidth}{0.5pt}\end{center}

\subsection{4.1 Tren Eskalasi Konflik Agraria Seiring Ekspansi
Industri}\label{tren-eskalasi-konflik-agraria-seiring-ekspansi-industri}

\textbf{Metode: Analisis Tren Time-Series (Sumber: KPA / Tanah Kita)}

\subsubsection{Metodologi: Analisis Tren
Time-Series}\label{metodologi-analisis-tren-time-series}

\textbf{Metode Analisis:} Sub-bab ini menggunakan visualisasi tren
runtun waktu (\emph{Time-Series Trend Analysis}) untuk melacak eskalasi
kasus perampasan lahan secara historis.

\begin{enumerate}
\def\labelenumi{\arabic{enumi}.}
\tightlist
\item
  \textbf{Model Analisis Tren Historis:}

  \begin{itemize}
  \tightlist
  \item
    \textbf{Time-Series Tracking:} Memetakan fluktuasi dan eskalasi
    frekuensi letupan konflik agraria dalam rentang waktu memanjang
    (longitudinal).
  \item
    \textbf{Komparasi Periodik:} Membandingkan volume letupan konflik
    antara fase pra-ekspansi (sebelum hilirisasi masif) dengan fase
    pasca-ekspansi (era Proyek Strategis Nasional).
  \item
    \textbf{Pemetaan Eskalasi:} Mengidentifikasi pola lonjakan kasus
    perampasan lahan untuk membuktikan secara empiris relasi antara
    percepatan industrialisasi dengan peningkatan konflik sosial.
  \end{itemize}
\item
  \textbf{Kalkulasi/Formula Pengolahan:} Agregasi jumlah konflik
  berdasarkan periode tahun pencatatan dan sektor industri.

  \begin{itemize}
  \tightlist
  \item
    \texttt{Total\_Konflik\_Tahunan\ =\ COUNT(Kasus)\ GROUP\ BY\ Tahun,\ Sektor}
  \item
    \texttt{Lonjakan\_Eskalasi\ =\ (Kasus\_Pasca\ -\ Kasus\_Pra)\ /\ Kasus\_Pra\ *\ 100\%}
  \end{itemize}
\item
  \textbf{Variabel \& Fitur Data:}

  \begin{itemize}
  \tightlist
  \item
    \textbf{Waktu (Independen):} Tahun pencatatan konflik (1990 - 2025).
  \item
    \textbf{Frekuensi \& Sektor (Dependen):} Jumlah insiden perampasan
    ruang dan sektor korporasi yang memicu konflik.
  \end{itemize}
\item
  \textbf{Dataset \& File:}

  \begin{itemize}
  \tightlist
  \item
    Catatan Konflik Agraria:
    \texttt{data/processed/sulawesi\_konflik\_agraria\_tanahkita.csv}
  \end{itemize}
\end{enumerate}

\begin{center}\rule{0.5\linewidth}{0.5pt}\end{center}

Visualisasi \emph{time-series} di bawah ini memberikan bukti empiris
yang tidak dapat dibantah mengenai korelasi langsung antara ekspansi
industri berskala masif dengan eskalasi letupan konflik agraria di
daratan Sulawesi. Secara historis, jika kita membandingkan dua periode
waktu yang berbeda, lonjakan perampasan ruang hidup masyarakat terlihat
sangat drastis dan tidak proporsional. Pada periode pra-2005, sistem
pendataan mencatat ``hanya'' terdapat \textbf{13 kasus} letupan konflik
yang tereskalasi. Angka ini secara fundamental merepresentasikan
dinamika agraria tradisional sebelum keran perizinan konsesi ekstraktif
dibuka secara agresif oleh pemerintah daerah pasca implementasi otonomi
daerah secara penuh.

Namun, narasi harmoni pembangunan ini hancur berantakan ketika memasuki
periode pasca-2005 hingga saat ini. Data empiris secara mengejutkan
mencatat setidaknya \textbf{82 kasus} perampasan lahan yang memicu
perlawanan berdarah, yang ekuivalen dengan lonjakan eskalasi raksasa
sebesar \textbf{630.8\%} dibandingkan era sebelumnya. Transformasi tata
ruang yang sangat brutal ini didorong oleh lahirnya rezim komodifikasi
daratan, di mana penerbitan Izin Usaha Pertambangan (IUP) mineral dan
batubara, serta ekspansi Hak Guna Usaha (HGU) untuk perkebunan kelapa
sawit monokultur menjadi panglima pembangunan yang menggusur wilayah
kelola masyarakat adat dan petani gurem. Hal ini secara faktual
membuktikan bahwa model pembangunan berorientasi PDB (Produk Domestik
Bruto) nyatanya beroperasi di atas kerentanan ruang hidup warga.

Lebih jauh lagi, jika membedah tren pada satu dekade terakhir (terutama
puncak eskalasi masif pada tahun 2017 dan melesat pasca-2020), kita
menemukan anomali yang sangat berbahaya. Tren letupan sengketa sosial
ini tidak lagi sekadar didominasi oleh perambahan hutan lindung atau
perluasan kebun sawit, melainkan telah bermutasi menjadi konflik
struktural akibat narasi besar \textbf{Hilirisasi Nikel} dan pengadaan
daratan secara darurat untuk \textbf{Proyek Strategis Nasional
(Infrastruktur \& PSN)}. Warga lokal dipaksa melepaskan hak atas tanah
produktif mereka di wilayah-wilayah episentrum ekstraktif demi menggelar
karpet merah bagi modal korporat transnasional. Fakta keras berupa
\textbf{95 total insiden historis} ini secara definitif membantah klaim
negara bahwa industrialisasi ekstraktif membawa efek kesejahteraan
berganda (\emph{trickle-down effect}). Sebaliknya, kawasan-kawasan
investasi tersebut justru bermetamorfosis menjadi `zona tumbal'
(\emph{sacrifice zones}) di mana laju akumulasi kapital segelintir elit
korporasi harus dibayar sangat mahal dengan ongkos krisis ekologis
permanen, represi aparat negara, serta hancurnya tatanan kedaulatan
pangan maupun pranata sosial masyarakat lokal.

\begin{figure}
\centering
\pandocbounded{\includegraphics[keepaspectratio,alt={Ledakan Konflik Agraria di Sulawesi (1990-2025)}]{visuals_bab4/chart_4_1_konflik_timeseries.png}}
\caption{Ledakan Konflik Agraria di Sulawesi (1990-2025)}
\end{figure}

\begin{quote}
\textbf{Interpretasi Ekologis: Anatomi Ledakan Konflik 2017}

Grafik di atas secara gamblang memperlihatkan anomali eskalasi ekstrem
yang memuncak pada \textbf{tahun 2017} dengan rekor \textbf{75 letupan
konflik}. Pembedahan data sektoral membongkar bahwa krisis ini bukanlah
sekadar kebetulan; ledakan ini didominasi secara mutlak oleh sektor
\textbf{Kehutanan (40 kasus)} dan \textbf{Perkebunan (21 kasus)}, yang
kemudian diikuti oleh penetrasi \textbf{Pertambangan dan Infrastruktur
PSN}. Tahun 2017 menandai periode kelam \emph{(inflection point)} di
mana pemerintah mengakselerasi pelepasan kawasan hutan dan Izin Pinjam
Pakai Kawasan Hutan (IPPKH) secara masif guna memfasilitasi rantai pasok
nikel dan megaproyek strategis nasional. Ekspansi spasial yang brutal
ini secara langsung merampas wilayah kelola masyarakat adat dan merusak
ekosistem penyangga, memicu gelombang perlawanan akar rumput yang
direpresi. Secara empiris, narasi hilirisasi telah membuktikan dirinya
beroperasi di atas ongkos perampasan ruang hidup berskala masif.
\end{quote}

\begin{quote}
\textbf{Interpretasi Ekologis dan Sosial:} Loncatan drastis letupan
konflik terjadi beririsan dengan agresivitas rezim perizinan. Hilirisasi
Nikel dan Proyek Strategis Nasional (PSN) secara faktual telah
merekayasa kawasan investasi menjadi zona tumbal yang mengorbankan
kedaulatan masyarakat lokal secara permanen.
\end{quote}

\begin{center}\rule{0.5\linewidth}{0.5pt}\end{center}

\subsection{4.2 Sebaran Sektoral: Korban Jiwa dan Monopoli
Ruang}\label{sebaran-sektoral-korban-jiwa-dan-monopoli-ruang}

\textbf{Metode: Analisis Komparatif Dampak Sosial-Ekologis (Sumber: KPA
/ Tanah Kita)}

\subsubsection{Metodologi: Analisis Komparatif Dampak
Sosial-Ekologis}\label{metodologi-analisis-komparatif-dampak-sosial-ekologis}

\textbf{Metode Analisis:} Sub-bab ini menggunakan agregasi komparatif
(\emph{Comparative Aggregation Analysis}) untuk membedah skala
kehancuran sosial (korban terdampak) dan monopoli ruang (hektar) antar
sektor.

\begin{enumerate}
\def\labelenumi{\arabic{enumi}.}
\tightlist
\item
  \textbf{Model Analisis Beban Sektoral (Sectoral Burden Analysis):}

  \begin{itemize}
  \tightlist
  \item
    \textbf{Kategorisasi Sektoral (Profiling):} Mengklasifikasikan
    sumber konflik (sektor Tambang, Perkebunan, Kehutanan, dll.) sebagai
    basis pengelompokan (\emph{grouping}).
  \item
    \textbf{Kuantifikasi Monopoli:} Menghitung total agregat luasan
    daratan (hektar) yang dirampas dan jumlah masyarakat (jiwa) yang
    terdampak per sektor industri.
  \item
    \textbf{Evaluasi Dominasi:} Membedah asimetri penguasaan ruang untuk
    mengidentifikasi sektor mana yang bertindak sebagai aktor dominan
    dalam praktik perampasan tanah (\emph{land grabbing}).
  \end{itemize}
\item
  \textbf{Kalkulasi/Formula Pengolahan:} Perhitungan sum/agregat dari
  seluruh korban jiwa (bukan korban meninggal, melainkan terdampak) dan
  hektar.

  \begin{itemize}
  \tightlist
  \item
    \texttt{Total\_Jiwa\_Terdampak\ =\ SUM(Jiwa)\ GROUP\ BY\ Sektor}
  \item
    \texttt{Total\_Monopoli\_Area\ =\ SUM(Hektar)\ GROUP\ BY\ Sektor}
  \end{itemize}
\item
  \textbf{Variabel \& Fitur Data:}

  \begin{itemize}
  \tightlist
  \item
    \textbf{Sektor (Independen):} Kategori proyek (Perkebunan,
    Kehutanan, Pertambangan, dll).
  \item
    \textbf{Korban Jiwa \& Luas Area (Dependen):} Jumlah orang terdampak
    (Jiwa) dan luas sengketa (Ha).
  \end{itemize}
\item
  \textbf{Dataset \& File:}

  \begin{itemize}
  \tightlist
  \item
    Dampak Konflik:
    \texttt{data/processed/sulawesi\_konflik\_agraria\_tanahkita.csv}
  \end{itemize}
\end{enumerate}

\begin{center}\rule{0.5\linewidth}{0.5pt}\end{center}

Konflik agraria bukanlah sebuah insiden terisolasi yang hanya berupa
sengketa batas tanah, melainkan instrumen sistematis dari akumulasi
modal yang beroperasi dengan menggusur paksa kehidupan manusia.
Visualisasi komparatif di bawah ini membongkar skala kehancuran sosial
dan ekologis yang diakibatkan oleh masing-masing sektor industri
ekstraktif. Ketika kita membedah total jumlah korban terdampak, data
menunjukkan realitas yang sangat mengerikan. \textbf{Sektor Kehutanan}
menjadi penyumbang terbesar krisis kemanusiaan dengan total korban
mencapai \textbf{21,886 jiwa}. Angka ini bukan sekadar statistik; ini
merepresentasikan masyarakat adat dan komunitas lokal yang ruang hidup
dan wilayah adatnya direnggut atas nama legalitas izin Hutan Tanaman
Industri (HTI) maupun klaim sepihak kawasan lindung oleh negara.

Menyusul di posisi kedua adalah \textbf{Sektor Pertambangan} yang telah
memakan korban sebanyak \textbf{54,658 jiwa}. Lonjakan korban di sektor
ini berhubungan langsung dengan ambisi hilirisasi mineral kritis
(terutama nikel) yang memaksa warga pesisir dan petani untuk melepaskan
ruang produksi mereka demi fasilitas \emph{smelter} dan pertambangan
terbuka. Masyarakat yang melawan seringkali dihadapkan pada represi
berlapis, mulai dari intimidasi preman korporasi hingga kriminalisasi
oleh aparat keamanan negara yang bertindak sebagai penjaga gawang
investasi.

Di sisi lain, saat kita meninjau dari dimensi monopoli tata ruang
(luasan hektar yang dikonflikkan), \textbf{Sektor
Perkebunan}---khususnya ekspansi kelapa sawit---menjadi penguasa absolut
dengan merampas lahan seluas \textbf{77,902 Hektar}. Konsentrasi
penguasaan tanah oleh segelintir korporasi perkebunan ini menghancurkan
kedaulatan pangan lokal dan menciptakan ketimpangan agraria yang
struktural. Disusul oleh sektor Kehutanan seluas \textbf{66,193 Ha} dan
Pertambangan seluas \textbf{441,286 Ha}, trinitas sektor ekstraktif ini
(Kebun, Hutan, Tambang) secara empiris membuktikan bahwa pembangunan
ekonomi selama ini semata-mata bergantung pada perampasan ruang berskala
masif. Tidak ada tetesan kesejahteraan (\emph{trickle-down effect}) bagi
warga tapak; yang tersisa hanyalah kemiskinan struktural, pencemaran
tanah, dan hilangnya hak-hak dasar konstitusional mereka atas daratan
yang telah mereka tempati secara turun-temurun.

\begin{figure}
\centering
\pandocbounded{\includegraphics[keepaspectratio,alt={Ledakan Korban Terdampak (Jiwa) per Tahun}]{visuals_bab4/chart_4_2a_jiwa.png}}
\caption{Ledakan Korban Terdampak (Jiwa) per Tahun}
\end{figure}

\begin{figure}
\centering
\pandocbounded{\includegraphics[keepaspectratio,alt={Monopoli Area Konflik (Hektar) per Tahun}]{visuals_bab4/chart_4_2b_ha.png}}
\caption{Monopoli Area Konflik (Hektar) per Tahun}
\end{figure}

\begin{quote}
\textbf{Interpretasi Ekologis dan Sosial:} Lonjakan luar biasa pada
grafik merepresentasikan titik didih ledakan demografis dari kegagalan
mutlak sistem pengaman sosial di zona investasi ekstraktif.
\end{quote}

\subsubsection{Bedah Forensik Anomali (Spike) Konflik
Agraria}\label{bedah-forensik-anomali-spike-konflik-agraria}

Berdasarkan ekstraksi dataset secara mendalam, berikut adalah bedah
anatomis dari lonjakan-lonjakan ekstrem (\emph{spikes}) yang terjadi
pada grafik \textbf{Ledakan Korban Terdampak (Jiwa)} dan
\textbf{Monopoli Area Konflik (Hektar)} di wilayah ini.

\begin{center}\rule{0.5\linewidth}{0.5pt}\end{center}

\subsection{4.3 Kriminalisasi Aktivis dan Resistensi Ruang
Sipil}\label{kriminalisasi-aktivis-dan-resistensi-ruang-sipil}

\textbf{Metode: Analisis Agregat Kasus Represi \& Pelanggaran HAM
(Sumber: Database Tanah Kita)}

\subsubsection{Metodologi: Analisis Agregat Kasus Represi \& Pelanggaran
HAM}\label{metodologi-analisis-agregat-kasus-represi-pelanggaran-ham}

\textbf{Metode Analisis:} Sub-bab ini menggunakan agregasi kasus
indikasi pelanggaran Hak Asasi Manusia dan Kriminalisasi Pejuang
Lingkungan melalui ekstraksi metrik fatalitas.

\begin{enumerate}
\def\labelenumi{\arabic{enumi}.}
\tightlist
\item
  \textbf{Pemodelan Indikator Kekerasan \& Represi:}

  \begin{itemize}
  \tightlist
  \item
    \textbf{Violence \& Criminalization Tracking:} Mendokumentasikan
    kasus penangkapan, intimidasi, kekerasan fisik, hingga jatuhnya
    korban jiwa di pihak warga dan aktivis lingkungan.
  \item
    \textbf{Kuantifikasi Fatalitas:} Menghitung akumulasi jumlah korban
    kriminalisasi dan korban tewas sebagai proksi tingkat represi
    struktural.
  \item
    \textbf{Pemetaan Ruang Sipil:} Mengevaluasi sejauh mana ekspansi
    investasi industri ekstraktif beroperasi dengan menggunakan
    instrumen represi aparatur keamanan (penyempitan ruang sipil).
  \end{itemize}
\item
  \textbf{Kalkulasi/Formula Pengolahan:} Penghitungan jumlah insiden
  kriminalisasi serta total akumulasi korban represi kekerasan fisik.

  \begin{itemize}
  \tightlist
  \item
    \texttt{Total\_Kasus\_Kriminalisasi\ =\ COUNT(Kasus)\ WHERE\ Indikasi\_Kriminalisasi\ =\ TRUE}
  \item
    \texttt{Total\_Korban\_Tewas\ =\ SUM(Jumlah\_Tewas)\ GROUP\ BY\ Sektor}
  \end{itemize}
\item
  \textbf{Variabel \& Fitur Data:}

  \begin{itemize}
  \tightlist
  \item
    \textbf{Status Represi (Dependen):} Boolean (Ya/Tidak) terjadinya
    indikasi kriminalisasi dalam konflik.
  \item
    \textbf{Kuantitas Korban (Dependen):} Angka mutlak (integer) korban
    tertangkap, terluka, dan meninggal.
  \end{itemize}
\item
  \textbf{Dataset \& File:}

  \begin{itemize}
  \tightlist
  \item
    Represi dan Kriminalisasi:
    \texttt{data/processed/sulawesi\_konflik\_agraria\_tanahkita.csv}
  \end{itemize}
\end{enumerate}

\begin{center}\rule{0.5\linewidth}{0.5pt}\end{center}

Rentetan data kuantitatif di wilayah Sulawesi secara telanjang membantah
klaim arus utama yang kerap didengungkan oleh pemerintah dan oligarki
korporasi, bahwa ekspansi industri ekstraktif membawa kesejahteraan dan
pertumbuhan inklusif bagi masyarakat lokal. Fakta empiris justru
memperlihatkan bahwa tata kelola investasi di Indonesia secara
struktural dibangun di atas fondasi represi dan kekerasan terhadap ruang
sipil.

Dari \textbf{46 kasus indikasi kriminalisasi} yang berhasil
didokumentasikan, tercatat sebanyak \textbf{93 warga dan aktivis
lingkungan yang ditangkap} secara sewenang-wenang. Angka ini bukanlah
statistik hampa, melainkan representasi dari hancurnya keadilan ekologis
dan perampasan ruang hidup masyarakat adat, petani, dan nelayan yang
dipaksa menyerahkan tanah leluhurnya demi akumulasi kapital segelintir
elit industri ekstraktif.

Jika kita membedah lebih dalam pada distribusi sektoral, \textbf{Sektor
Pertambangan} muncul sebagai aktor dominan yang paling sering
menggunakan instrumen koersif negara, menyumbang total \textbf{17 kasus
represi}. Penggunaan aparat keamanan negara maupun preman korporasi
untuk memuluskan perampasan tanah menunjukkan bahwa hukum seringkali
ditundukkan pada kepentingan bisnis raksasa yang lapar lahan. Eskalasi
konflik paling mematikan mencapai puncaknya pada tahun \textbf{2017}
dengan mencatatkan \textbf{7 kasus secara bersamaan}. Dalam banyak
peristiwa empiris, warga lokal yang sekadar mempertahankan hak
konstitusional mereka atas lingkungan hidup yang baik dan sehat justru
dilabeli sebagai provokator dan dijerat pasal pidana karet.

Tragedi kemanusiaan ini menjadi semakin kelam dengan hilangnya nyawa
\textbf{1 pejuang lingkungan} yang melayang sia-sia di pusaran konflik
agraria. Gugurnya pahlawan-pahlawan ruang hidup ini menggarisbawahi
kegagalan mutlak instrumen pengaman ekologis - seperti D3TLH maupun
dokumen AMDAL - dalam menjamin keselamatan rakyat. Selama pendekatan
pembangunan eksploitatif yang bertumpu pada sekuritisasi investasi ini
dipertahankan, setiap hektar hutan yang dibabat akan selalu berlumuran
air mata konflik.

{\def\LTcaptype{none} % do not increment counter
\begin{longtable}[]{@{}
  >{\raggedright\arraybackslash}p{(\linewidth - 6\tabcolsep) * \real{0.2500}}
  >{\raggedright\arraybackslash}p{(\linewidth - 6\tabcolsep) * \real{0.2500}}
  >{\raggedright\arraybackslash}p{(\linewidth - 6\tabcolsep) * \real{0.2500}}
  >{\raggedright\arraybackslash}p{(\linewidth - 6\tabcolsep) * \real{0.2500}}@{}}
\toprule\noalign{}
\begin{minipage}[b]{\linewidth}\raggedright
Kasus Indikasi Kriminalisasi
\end{minipage} & \begin{minipage}[b]{\linewidth}\raggedright
Warga/Aktivis Ditangkap
\end{minipage} & \begin{minipage}[b]{\linewidth}\raggedright
Korban Luka-luka
\end{minipage} & \begin{minipage}[b]{\linewidth}\raggedright
Korban Tewas
\end{minipage} \\
\midrule\noalign{}
\endhead
\bottomrule\noalign{}
\endlastfoot
\textbf{46 Kasus} & \textbf{93 Orang} & \textbf{4 Orang} & \textbf{1
Orang} \\
\end{longtable}
}

\begin{figure}
\centering
\pandocbounded{\includegraphics[keepaspectratio,alt={Tren Kasus Kriminalisasi \& Represi (Pasca 2000)}]{visuals_bab4/chart_4_3a_kriminalisasi_trend.png}}
\caption{Tren Kasus Kriminalisasi \& Represi (Pasca 2000)}
\end{figure}

\begin{figure}
\centering
\pandocbounded{\includegraphics[keepaspectratio,alt={Sektor Industri Paling Represif}]{visuals_bab4/chart_4_3b_sektor_represif.png}}
\caption{Sektor Industri Paling Represif}
\end{figure}

\begin{quote}
\textbf{Interpretasi Ekologis \& Hak Asasi Manusia:} Tingginya angka
kriminalisasi dan korban tewas di sekitar area konsesi (terutama
Pertambangan) membuktikan bahwa perampasan ruang selalu dibarengi dengan
pendekatan represif. Ini membantah telak narasi ``Hilirisasi Hijau''
yang nyatanya ditebus dengan ongkos kemanusiaan yang berdarah.
\end{quote}

\paragraph{Arsip Kasus Represi dan Kekerasan Fisik
Tertinggi}\label{arsip-kasus-represi-dan-kekerasan-fisik-tertinggi}

\emph{Menampilkan 10 kasus dengan jumlah korban penangkapan atau tewas
terbanyak berdasarkan data yang berhasil didokumentasikan.}

\begin{center}\rule{0.5\linewidth}{0.5pt}\end{center}

\subsection{4.4 Pembuktian Statistik: Ekspansi vs Eskalasi
Konflik}\label{pembuktian-statistik-ekspansi-vs-eskalasi-konflik}

\textbf{Metode: Before-After Analysis \& Crosstabulation}

\subsubsection{Metodologi: Before-After Analysis \&
Crosstabulation}\label{metodologi-before-after-analysis-crosstabulation}

\textbf{Metode Analisis:} Sub-bab ini menggunakan Uji Chi-Square
(\emph{Crosstabulation}) dan kalkulasi risiko peluang (\emph{Odds
Ratio}) untuk menguji validitas empiris secara akademis.

\begin{enumerate}
\def\labelenumi{\arabic{enumi}.}
\tightlist
\item
  \textbf{Uji Korelasi Variabel Kategorikal:}

  \begin{itemize}
  \tightlist
  \item
    \textbf{Crosstabulation:} Mentabulasi silang frekuensi kemunculan
    dua kondisi (Contoh: Keterlibatan Perusahaan vs Adanya
    Kriminalisasi) untuk mencari relasi ketergantungan.
  \item
    \texttt{H0\ (Null\ Hypothesis):\ Variabel\ baris\ (Periode/Aktor)\ saling\ bebas\ (independent)\ secara\ absolut\ terhadap\ variabel\ kolom\ (Represi/Kematian).}
  \item
    \texttt{Decision\ Rule:\ Chi-Square\ Asymptotic\ Significance\ (P-Value)\ \textless{}\ 0.05,\ maka\ tolak\ H0\ (Terdapat\ korelasi\ yang\ signifikan).}
  \end{itemize}
\item
  \textbf{Kalkulasi/Formula Pengolahan:} Algoritma Uji Tabulasi Silang
  Chi-Square.

  \begin{itemize}
  \tightlist
  \item
    \texttt{Chi-Square\ (χ²)\ =\ Σ\ {[}(Observed\ -\ Expected)²\ /\ Expected{]}}
  \item
    \texttt{Odds\ Ratio\ (OR)\ =\ (Sel\ A\ ×\ Sel\ D)\ /\ (Sel\ B\ ×\ Sel\ C)}
  \end{itemize}
\item
  \textbf{Variabel \& Fitur Data:}

  \begin{itemize}
  \tightlist
  \item
    \textbf{Matriks Ekspansi (Independen):} Dikotomi rentang waktu
    (Pra/Pasca 2014) dan kehadiran korporasi.
  \item
    \textbf{Matriks Eskalasi (Dependen):} Kehadiran status represi dan
    terjadinya jatuhnya korban nyawa (Boolean dikonversi ke kategori).
  \end{itemize}
\item
  \textbf{Dataset \& File:}

  \begin{itemize}
  \tightlist
  \item
    Base Data Cross-Section:
    \texttt{data/processed/sulawesi\_konflik\_agraria\_tanahkita.csv}
  \end{itemize}
\end{enumerate}

\begin{center}\rule{0.5\linewidth}{0.5pt}\end{center}

Hipotesis utama dalam evaluasi ini adalah bahwa \textbf{industrialisasi
dan ekspansi korporasi} berbanding lurus dengan \textbf{eskalasi konflik
dan represi} terhadap masyarakat. Untuk mengujinya secara statistik
sesuai pedoman D3TLH, analisis dibagi menjadi dua bagian: (1) Komparasi
metrik Before-After, dan (2) Uji signifikansi Crosstab Chi-Square. Unit
observasinya adalah catatan kejadian letupan konflik historis.

\subsubsection{A. Analisis Komparatif Before-After (Pra vs Era
Hilirisasi)}\label{a.-analisis-komparatif-before-after-pra-vs-era-hilirisasi}

Perbandingan absolut eskalasi konflik agraria sebelum dan sesudah rezim
hilirisasi masif dimulai (cut-off tahun 2014).

{\def\LTcaptype{none} % do not increment counter
\begin{longtable}[]{@{}
  >{\raggedright\arraybackslash}p{(\linewidth - 8\tabcolsep) * \real{0.2000}}
  >{\raggedright\arraybackslash}p{(\linewidth - 8\tabcolsep) * \real{0.2000}}
  >{\raggedright\arraybackslash}p{(\linewidth - 8\tabcolsep) * \real{0.2000}}
  >{\raggedright\arraybackslash}p{(\linewidth - 8\tabcolsep) * \real{0.2000}}
  >{\raggedright\arraybackslash}p{(\linewidth - 8\tabcolsep) * \real{0.2000}}@{}}
\toprule\noalign{}
\begin{minipage}[b]{\linewidth}\raggedright
Periode
\end{minipage} & \begin{minipage}[b]{\linewidth}\raggedright
Rata-rata Konflik
\end{minipage} & \begin{minipage}[b]{\linewidth}\raggedright
Total Letupan
\end{minipage} & \begin{minipage}[b]{\linewidth}\raggedright
Warga Ditangkap
\end{minipage} & \begin{minipage}[b]{\linewidth}\raggedright
Korban Tewas
\end{minipage} \\
\midrule\noalign{}
\endhead
\bottomrule\noalign{}
\endlastfoot
\textbf{Pra-Ekspansi (1990 -- 2013)} & \textbf{2.1 Kasus/Tahun} & 37
kejadian & 83 jiwa & 0 jiwa \\
\textbf{Pasca-Ekspansi (2014 -- 2024)} & \textbf{5.5 Kasus/Tahun} & 55
kejadian & 10 jiwa & 1 jiwa \\
\end{longtable}
}

\subsubsection{B. Uji Statistik Crosstab
(Chi-Square)}\label{b.-uji-statistik-crosstab-chi-square}

\textbf{Variabel Independen (X):} Periode Ekspansi Industri

\textbf{Variabel Dependen (Y):} Tingkat Represi \& Kriminalisasi

\paragraph{Case Processing Summary}\label{case-processing-summary}

{\def\LTcaptype{none} % do not increment counter
\begin{longtable}[]{@{}
  >{\raggedright\arraybackslash}p{(\linewidth - 12\tabcolsep) * \real{0.1429}}
  >{\raggedright\arraybackslash}p{(\linewidth - 12\tabcolsep) * \real{0.1429}}
  >{\raggedright\arraybackslash}p{(\linewidth - 12\tabcolsep) * \real{0.1429}}
  >{\raggedright\arraybackslash}p{(\linewidth - 12\tabcolsep) * \real{0.1429}}
  >{\raggedright\arraybackslash}p{(\linewidth - 12\tabcolsep) * \real{0.1429}}
  >{\raggedright\arraybackslash}p{(\linewidth - 12\tabcolsep) * \real{0.1429}}
  >{\raggedright\arraybackslash}p{(\linewidth - 12\tabcolsep) * \real{0.1429}}@{}}
\toprule\noalign{}
\begin{minipage}[b]{\linewidth}\raggedright
\end{minipage} & \begin{minipage}[b]{\linewidth}\raggedright
Valid N
\end{minipage} & \begin{minipage}[b]{\linewidth}\raggedright
Valid \%
\end{minipage} & \begin{minipage}[b]{\linewidth}\raggedright
Missing N
\end{minipage} & \begin{minipage}[b]{\linewidth}\raggedright
Missing \%
\end{minipage} & \begin{minipage}[b]{\linewidth}\raggedright
Total N
\end{minipage} & \begin{minipage}[b]{\linewidth}\raggedright
Total \%
\end{minipage} \\
\midrule\noalign{}
\endhead
\bottomrule\noalign{}
\endlastfoot
Periode Ekspansi Industri * Tingkat Represi \& Kriminalisasi & 523 &
100.0\% & 0 & 0.0\% & 523 & 100.0\% \\
\end{longtable}
}

\paragraph{Periode Ekspansi Industri * Tingkat Represi \& Kriminalisasi
Crosstabulation}\label{periode-ekspansi-industri-tingkat-represi-kriminalisasi-crosstabulation}

{\def\LTcaptype{none} % do not increment counter
\begin{longtable}[]{@{}llll@{}}
\toprule\noalign{}
& Baseline (Tanpa Kriminalisasi) & Ada Represi/Kriminalisasi & Total \\
\midrule\noalign{}
\endhead
\bottomrule\noalign{}
\endlastfoot
\textbf{Pra-ekspansi (\textless{} 2014)} Count & 179 & 73 & 252 \\
\textbf{Pra-ekspansi (\textless{} 2014)} Expected & 157.6 & 94.4 &
252.0 \\
\textbf{Pasca-ekspansi (≥ 2014)} Count & 148 & 123 & 271 \\
\textbf{Pasca-ekspansi (≥ 2014)} Expected & 169.4 & 101.6 & 271.0 \\
\textbf{Total} Count & 327 & 196 & 523 \\
\textbf{Total} Expected & 327.0 & 196.0 & 523.0 \\
\end{longtable}
}

\paragraph{Chi-Square Tests}\label{chi-square-tests}

\textbf{Periode Ekspansi Industri * Tingkat Represi \& Kriminalisasi}

{\def\LTcaptype{none} % do not increment counter
\begin{longtable}[]{@{}llll@{}}
\toprule\noalign{}
& Value & df & Asymp. Sig. (2-sided) \\
\midrule\noalign{}
\endhead
\bottomrule\noalign{}
\endlastfoot
Pearson Chi-Square & 14.331 & 1 & 0.000 \\
Likelihood Ratio & 14.447 & 1 & 0.000 \\
Linear-by-Linear Association & 14.995 & 1 & 0.000 \\
N of Valid Cases & 523 & & \\
\end{longtable}
}

\subsubsection{Ringkasan Uji Hipotesis}\label{ringkasan-uji-hipotesis}

\textbf{Result: SIGNIFIKAN (Ada Hubungan)}

{\def\LTcaptype{none} % do not increment counter
\begin{longtable}[]{@{}ll@{}}
\toprule\noalign{}
Parameter & Nilai \\
\midrule\noalign{}
\endhead
\bottomrule\noalign{}
\endlastfoot
P-Value & 0.0002 \\
Chi-Square & 14.331 \\
df & 1 \\
\textbf{Odds Ratio (Risk Estimate)} & \textbf{2.038} \\
\end{longtable}
}

\begin{quote}
\textbf{Interpretasi Sosial Kritis:} Temuan ini sangat krusial:
pergeseran status \textbf{Periode Ekspansi Industri} terbukti
\textbf{berkorelasi kuat dan signifikan} dengan \textbf{Tingkat Represi
\& Kriminalisasi} (P \textless{} 0.05). Angka Odds Ratio (OR: 2.038)
menjadi konfirmasi empiris bahwa narasi hilirisasi dan investasi
bukanlah agenda nirkekerasan---ekspansi spasial mereka mutlak
mengeskalasi pelanggaran hak asasi masyarakat tapak.
\end{quote}

\begin{center}\rule{0.5\linewidth}{0.5pt}\end{center}

\subsubsection{Ringkasan Eksekutif Seluruh Skenario
Crosstab}\label{ringkasan-eksekutif-seluruh-skenario-crosstab}

Tabel di bawah ini merangkum hasil pengujian statistik (Chi-Square)
untuk semua kemungkinan kombinasi antara indikator Ekspansi (X) dan
Eskalasi Konflik (Y) pada panel data yang sama.

{\def\LTcaptype{none} % do not increment counter
\begin{longtable}[]{@{}
  >{\raggedright\arraybackslash}p{(\linewidth - 10\tabcolsep) * \real{0.1667}}
  >{\raggedright\arraybackslash}p{(\linewidth - 10\tabcolsep) * \real{0.1667}}
  >{\raggedright\arraybackslash}p{(\linewidth - 10\tabcolsep) * \real{0.1667}}
  >{\raggedright\arraybackslash}p{(\linewidth - 10\tabcolsep) * \real{0.1667}}
  >{\raggedright\arraybackslash}p{(\linewidth - 10\tabcolsep) * \real{0.1667}}
  >{\raggedright\arraybackslash}p{(\linewidth - 10\tabcolsep) * \real{0.1667}}@{}}
\toprule\noalign{}
\begin{minipage}[b]{\linewidth}\raggedright
Variabel Independen (X)
\end{minipage} & \begin{minipage}[b]{\linewidth}\raggedright
Variabel Dependen (Y)
\end{minipage} & \begin{minipage}[b]{\linewidth}\raggedright
Chi-Square
\end{minipage} & \begin{minipage}[b]{\linewidth}\raggedright
P-Value
\end{minipage} & \begin{minipage}[b]{\linewidth}\raggedright
Odds Ratio
\end{minipage} & \begin{minipage}[b]{\linewidth}\raggedright
Kesimpulan
\end{minipage} \\
\midrule\noalign{}
\endhead
\bottomrule\noalign{}
\endlastfoot
Periode Ekspansi Industri & Tingkat Represi \& Kriminalisasi & 14.331 &
0.000 & 2.04 & 🟢 SIGNIFIKAN \\
Periode Ekspansi Industri & Tingkat Penelantaran Kasus & 3.283 & 0.070 &
1.40 & 🔴 TIDAK SIGNIFIKAN \\
Periode Ekspansi Industri & Tingkat Insiden Fisik (Luka/Tewas/Ditangkap)
& 2.636 & 0.104 & 0.53 & 🔴 TIDAK SIGNIFIKAN \\
Tipe Sektor (Tambang vs Non-Tambang) & Tingkat Represi \& Kriminalisasi
& 13.609 & 0.000 & 2.84 & 🟢 SIGNIFIKAN \\
Tipe Sektor (Tambang vs Non-Tambang) & Tingkat Penelantaran Kasus &
4.742 & 0.029 & 1.89 & 🟢 SIGNIFIKAN \\
Tipe Sektor (Tambang vs Non-Tambang) & Tingkat Insiden Fisik
(Luka/Tewas/Ditangkap) & 0.665 & 0.415 & 1.66 & 🔴 TIDAK SIGNIFIKAN \\
Keterlibatan Aparat/Pemerintah & Tingkat Represi \& Kriminalisasi &
55.633 & 0.000 & 4.80 & 🟢 SIGNIFIKAN \\
Keterlibatan Aparat/Pemerintah & Tingkat Penelantaran Kasus & 1.135 &
0.287 & 0.81 & 🔴 TIDAK SIGNIFIKAN \\
Keterlibatan Aparat/Pemerintah & Tingkat Insiden Fisik
(Luka/Tewas/Ditangkap) & 3.253 & 0.071 & 2.23 & 🔴 TIDAK SIGNIFIKAN \\
\end{longtable}
}

\begin{quote}
\textbf{Pembedahan Realitas Kemanusiaan:}

Dari \textbf{9 skenario pengujian}, terdapat \textbf{4 skenario yang
terbukti SIGNIFIKAN}.
\end{quote}

Angka-angka pada tabel di atas bukan sekadar statistik di atas kertas,
melainkan \textbf{bukti empiris} dari brutalitas pembangunan. Tingginya
angka kemunculan represi pada skenario yang signifikan menegaskan bahwa
setiap kali wilayah operasi investasi diperlebar, probabilitas
dihadapkannya moncong senjata kepada warga melonjak drastis.

Skenario yang \emph{TIDAK SIGNIFIKAN} tidak berarti rezim terbebas dari
dosa kekerasan, melainkan bukti bahwa represi terhadap warga yang
mempertahankan tanahnya telah menjadi kultur mapan yang menyebar secara
sporadis melampaui sekat waktu dan korporasi.

\begin{center}\rule{0.5\linewidth}{0.5pt}\end{center}

\subsection{4.5 Peta Orkestrasi Konflik: Aktor Sipil vs Aktor
Ekstraktif}\label{peta-orkestrasi-konflik-aktor-sipil-vs-aktor-ekstraktif}

\textbf{Metode: Frequency Profiling (Text Parsing NLP) pada Data
TanahKita}

\subsubsection{Metodologi: Frequency Profiling (Text Parsing
NLP)}\label{metodologi-frequency-profiling-text-parsing-nlp}

\textbf{Metode Analisis:} Sub-bab ini menggunakan teknik pemrosesan teks
berbasis \emph{Natural Language Processing} (Regex Entity Extraction)
untuk membedah relasi aktor (korporasi vs sipil).

\begin{enumerate}
\def\labelenumi{\arabic{enumi}.}
\tightlist
\item
  \textbf{Model Ekstraksi Aktor (Entity Parsing \& Text Mining):}

  \begin{itemize}
  \tightlist
  \item
    \textbf{Textual Pattern Matching:} Memindai ribuan korpus teks
    narasi historis menggunakan metode \emph{Regular Expressions}
    (RegEx) untuk mendeteksi entitas korporasi (PT/CV) dan organisasi
    masyarakat sipil (CSO).
  \item
    \textbf{Token Counting (Frequency Profiling):} Menghitung frekuensi
    absolut penyebutan (\emph{mentions}) dari setiap aktor spesifik di
    dalam dokumentasi konflik.
  \item
    \textbf{Pemetaan Oligarki:} Memvalidasi indikasi konsentrasi
    kekuasaan dan monopoli penguasaan ruang oleh segelintir konglomerasi
    besar melalui seberapa sering nama entitas tersebut muncul dalam
    sengketa tanah.
  \end{itemize}
\item
  \textbf{Kalkulasi/Formula Pengolahan:} Regex pattern matching and
  Token Counting.

  \begin{itemize}
  \tightlist
  \item
    \texttt{Count\_PT\ =\ SUM(RegEx\_Match(r"\textbackslash{}b(?:PT\textbar{}CV)\textbackslash{}.?\textbackslash{}s*{[}A-Z{]}{[}a-zA-Z{]}*..."))}
  \item
    \texttt{Count\_CSO\ =\ SUM(RegEx\_Match(r"\textbackslash{}b(?:Walhi\textbar{}Jatam\textbar{}AMAN\textbar{}Aliansi)..."))}
  \end{itemize}
\item
  \textbf{Variabel \& Fitur Data:}

  \begin{itemize}
  \tightlist
  \item
    \textbf{Teks Korpus Historis (Independen):} Penggabungan kolom
    \texttt{judul}, \texttt{deskripsi}, dan \texttt{narasi} dari
    repositori kasus.
  \item
    \textbf{Frekuensi Penyebutan (Dependen):} \emph{Word counts}
    eksistensi entitas pada teks-teks sengketa.
  \end{itemize}
\item
  \textbf{Dataset \& File:}

  \begin{itemize}
  \tightlist
  \item
    Teks Bebas (\emph{Free-Text}):
    \texttt{data/processed/sulawesi\_konflik\_agraria\_tanahkita.csv}
  \end{itemize}
\end{enumerate}

\begin{center}\rule{0.5\linewidth}{0.5pt}\end{center}

Konflik yang membara tidak hanya melibatkan negara dan aparat, melainkan
memunculkan fenomena adu domba struktural (\emph{orkestrasi konflik
horizontal}). Pemecahan entitas (\emph{string parsing}) terhadap catatan
kronologi advokasi TanahKita menelanjangi siapa yang sesungguhnya
bermain di lapangan. Di satu sisi, masyarakat asli sering kali
didampingi oleh organisasi struktural yang solid, namun di sisi lain,
mulai muncul ormas-ormas, lembaga swadaya buatan, hingga institusi
pseudo-adat yang digunakan sebagai proksi (\emph{buffer}) oleh
korporasi. Grafik frekuensi ini membongkar dominasi aktor-aktor sipil
dan perusahaan tambang yang paling banyak merebut ruang hidup.

\paragraph{Top 10 Entitas Korporasi Paling
Dominan}\label{top-10-entitas-korporasi-paling-dominan}

\begin{figure}
\centering
\pandocbounded{\includegraphics[keepaspectratio,alt={Top 10 Entitas Korporasi Paling Dominan}]{visuals_bab4/chart_4_5a_korporasi.png}}
\caption{Top 10 Entitas Korporasi Paling Dominan}
\end{figure}

\begin{quote}
\textbf{Analisis Kritis:} Ekstraksi presisi tinggi membuktikan dominasi
absolut dari entitas \textbf{PT Perkebunan Nusantara (PTPN)} yang
terlibat dalam \textbf{153 catatan konflik terpisah}. Konsentrasi tinggi
frekuensi korporasi besar ini menegaskan bahwa represi di Sulawesi bukan
sekadar residu administratif, melainkan \emph{modus operandi} struktural
para penguasa modal skala masif.
\end{quote}

\paragraph{Top Aktor Proksi \& Vigilante
Terdeteksi}\label{top-aktor-proksi-vigilante-terdeteksi}

\begin{figure}
\centering
\pandocbounded{\includegraphics[keepaspectratio,alt={Top Aktor Proksi \& Vigilante Terdeteksi}]{visuals_bab4/chart_4_5b_vigilante.png}}
\caption{Top Aktor Proksi \& Vigilante Terdeteksi}
\end{figure}

\begin{quote}
\textbf{Analisis Kritis:} Kemunculan kelompok sipil seperti
\textbf{Preman} (terdeteksi hingga \textbf{16 kali}) menangkap besarnya
skala orkestrasi horizontal. Korporasi seringkali menggunakan jasa
pengamanan swakarsa, kelompok preman, hingga ormas vigilante sebagai
``bemper proksi'' untuk mengintimidasi warga lokal dan memecah belah
solidaritas akar rumput.
\end{quote}

\emph{* Grafik di atas hanya menampilkan Top 10 entitas. Untuk melihat
daftar lengkap dan detail seluruh aktor yang terdeteksi, silakan buka
tabel data di bawah ini.}

\emph{Sumber File:
\texttt{data/processed/sulawesi\_konflik\_agraria\_tanahkita.csv} - Data
diekstraksi secara dinamis menggunakan NLP Regex dari korpus narasi
seluruh kasus agraria (Nasional, N=568 kasus) untuk memetakan orkestrasi
struktural dan modus operandi aktor secara utuh.}

\end{document}
